\documentclass[11pt]{amsart}

\usepackage[T1]{fontenc}
\usepackage{lmodern}
\usepackage{microtype}
\usepackage{amsmath,amssymb,amsthm}
\usepackage{booktabs}
\usepackage[hidelinks]{hyperref}

\hypersetup{
  pdftitle={A Counterexample to the Lau--Shiu--Ng Conjecture on Moebius Ladders}
}

\newtheorem{theorem}{Theorem}
\theoremstyle{remark}
\newtheorem*{remark}{Remark}

\newcommand{\la}{\mathrm{la}}

\title[The Lau--Shiu--Ng conjecture on M\"obius ladders]
{A Counterexample to the Lau--Shiu--Ng Conjecture on M\"obius Ladders}

\subjclass[2020]{05C78, 05C15}
\keywords{local antimagic labeling, local antimagic chromatic number,
M\"obius ladder, Wagner graph}
\date{}

\begin{document}

\begin{abstract}
We give a local antimagic labeling of the eight-vertex M\"obius ladder
$M_8$ with three induced colors and prove that
$\chi_{\la}(M_8)=3$. This refutes Conjecture~3.14 of Lau, Shiu, and Ng
at its first admissible parameter. The same labeling also gives a
counterexample to Theorem~2.4 of arXiv:2203.06594v1, published as
Acta Math. Hungar. \textbf{169} (2023), 108--133.
\end{abstract}

\maketitle

\section{The counterexample}

Let $G$ be a graph with $q$ edges. A \emph{local antimagic labeling} is
a bijection
\[
f\colon E(G)\longrightarrow \{1,\dots,q\}
\]
such that the induced vertex sums
\[
f^+(x):=\sum_{e\ni x}f(e)
\]
satisfy $f^+(x)\ne f^+(y)$ whenever $xy\in E(G)$. The local antimagic
chromatic number $\chi_{\la}(G)$ is the minimum number of distinct
induced vertex sums over all local antimagic labelings of $G$.

For $n\ge2$, let $M_{2n}$ be obtained from the cycle
\[
u_1u_2\cdots u_nv_1v_2\cdots v_nu_1
\]
by adding the edges $u_iv_i$ for $1\le i\le n$. Lau, Shiu, and Ng
conjectured that $\chi_{\la}(M_{2n})=4$ for every even $n\ge4$
\cite[Conjecture~3.14]{LauShiuNg}. The first case covered by the
conjecture is $M_8$.

\begin{theorem}
The M\"obius ladder $M_8$ satisfies
\[
\chi_{\la}(M_8)=3.
\]
Consequently, Conjecture~3.14 of Lau, Shiu, and Ng is false.
\end{theorem}

\begin{proof}
Define $f\colon E(M_8)\to\{1,\dots,12\}$ by
\[
\begin{array}{c@{\,\mapsto\,}r@{\qquad}
                c@{\,\mapsto\,}r@{\qquad}
                c@{\,\mapsto\,}r}
\toprule
u_1u_2&1&u_2u_3&8&u_3u_4&3\\
u_4v_1&7&v_1v_2&5&v_2v_3&2\\
v_3v_4&11&v_4u_1&4&u_1v_1&12\\
u_2v_2&10&u_3v_3&6&u_4v_4&9\\
\bottomrule
\end{array}
\]
The labels are exactly $1,\dots,12$. Direct calculation gives
\[
\begin{array}{c|cccccccc}
x&u_1&u_2&u_3&u_4&v_1&v_2&v_3&v_4\\
\hline
f^+(x)&17&19&17&19&24&17&19&24.
\end{array}
\]
Thus the fibers of $f^+$ are
\[
\{u_1,u_3,v_2\},\qquad
\{u_2,u_4,v_3\},\qquad
\{v_1,v_4\},
\]
with induced colors $17$, $19$, and $24$, respectively. Each fiber is
an independent set in $M_8$, so $f$ is a local antimagic labeling and
$\chi_{\la}(M_8)\le3$.

The vertices $u_1,u_2,u_3,u_4,v_1$ form a $5$-cycle. Hence
$\chi(M_8)\ge3$. Since every local antimagic labeling induces a proper
vertex coloring,
\[
\chi_{\la}(M_8)\ge\chi(M_8)\ge3.
\]
Therefore $\chi_{\la}(M_8)=3$. Since $M_8=M_{2n}$ for the admissible
parameter $n=4$, the conjecture is false.
\end{proof}

\begin{remark}
Theorem~2.4 of arXiv:2203.06594v1 reads, in the notation used here,
``Suppose $G$ is a cubic graph of order $8$, then $\chi_{\la}(G)=4$''
\cite{LauShiuPreprint}. It stands in a subsection titled ``Connected
Cubic Graphs of Order at Most 8'', and its proof reduces the claim to
the four graphs $M_8,H_1,H_2,H_3$ enumerated there, so the intended
scope is the connected cubic graphs of order $8$. Since $M_8$ is a
connected cubic graph of order $8$ (indeed the first of those four
graphs), the theorem above also disproves that statement. That proof
uses
\[
2x+3y+3z=132
\]
for a hypothetical three-coloring with color-class sizes $2,3,3$.
For any labeling of a $12$-edge graph by $1,\dots,12$, however,
\[
\sum_{v\in V(G)}f^+(v)
 =2\sum_{e\in E(G)}f(e)
 =2\sum_{i=1}^{12}i
 =156.
\]
The value $132$ equals $2\sum_{i=1}^{11}i$, so the labels were summed as
if there were eleven of them rather than twelve. With the correct total,
the last of the five pairs $(x,\{y,z\})$ that the source derives from its
identity, namely $(12,\{17,19\})$, becomes $(24,\{17,19\})$, and this is
the pair realized by the labeling above:
\[
2\cdot24+3\cdot17+3\cdot19=156.
\]
The quotation and the identity above are taken from the v1 e-print; we
have not compared them with the corresponding text of the published
version.
The forward citations of \cite{LauShiuNg} and of arXiv:2203.06594 were
surveyed; the value $\chi_{\la}(M_8)=3$ does not appear among them, and
no erratum to either source is known to us.
No claim is made here about $\chi_{\la}(M_{2n})$ for even $n\ge6$.
\end{remark}

\begin{thebibliography}{9}

\bibitem{LauShiuNg}
G.-C. Lau, W.-C. Shiu, and H.-K. Ng,
\emph{On local antimagic chromatic number of cycle-related join graphs},
Discuss. Math. Graph Theory \textbf{41} (2021), 133--152.
\href{https://doi.org/10.7151/dmgt.2177}{doi:10.7151/dmgt.2177}.

\bibitem{LauShiuPreprint}
G.-C. Lau and W.-C. Shiu,
\emph{On join product and local antimagic chromatic number of regular graphs},
arXiv:2203.06594v1, 2022; published as
Acta Math. Hungar. \textbf{169} (2023), 108--133,
\href{https://doi.org/10.1007/s10474-023-01298-7}{doi:10.1007/s10474-023-01298-7}.

\end{thebibliography}

\end{document}